% appP_graph_c.tex -- generated by writing/tools/blueprint_gen.py 1.0 (level 3) on 2026-10-14 from sources c62e043e04; do not edit
\providecommand{\bpNum}[1]{\hyperref[#1]{\ref*{#1}}}%
\providecommand{\bpNumA}[1]{\hyperref[#1]{\textcolor{orange!65!black}{\ref*{#1}}}}%
\providecommand{\bpPage}[1]{\pageref*{#1}}%
\providecommand{\bpLab}[1]{{\ttfamily #1}}%
\providecommand{\bpBrk}{\allowbreak}%
\providecommand{\bpT}[2]{\facttagbox[#1]{\scriptsize #2}}%
\providecommand{\bpU}[1]{{\color{gray!70!black}#1}}%
\begin{figure}[tbp]
\centering
\begin{tikzpicture}[
  bp node/.style={draw,line width=0.45pt,rounded corners=1pt,font=\tiny\ttfamily,align=center,
                  inner xsep=1.5pt,inner ysep=1.1pt,minimum height=4mm},
  bp R/.style={bp node,fill=OliveGreen!15,draw=OliveGreen!60!black},
  bp A/.style={bp node,fill=RoyalBlue!15,draw=RoyalBlue!60!black},
  bp P/.style={bp node,fill=BurntOrange!15,draw=BurntOrange!70!black},
  bp C/.style={bp node,fill=Plum!15,draw=Plum!70!black},
  bp V/.style={bp node,fill=BrickRed!15,draw=BrickRed!70!black,double,double distance=0.6pt},
  bp cond/.style={dash pattern=on 1.6pt off 1pt},
  bp math/.style={bp node,fill=white,draw=black!60},
  bp def/.style={bp node,fill=white,draw=black!45,line width=0.3pt,rounded corners=0pt},
  bp conj/.style={bp node,ellipse,inner xsep=0pt,inner ysep=0pt,fill=Plum!15,draw=Plum!70!black},
  bp asm/.style={bp node,fill=orange!8,draw=orange!65!black,dash pattern=on 2pt off 1.2pt},
  bp stub/.style={bp node,font=\tiny,fill=gray!12,draw=gray!55,text=black!70},
  bp fold/.style={bp node,font=\tiny,fill=white,draw=gray!55,dash pattern=on 1pt off 1pt,text=black!70},
  bp unit/.style={font=\scriptsize,text=black!70,inner sep=1pt},
  bp sep/.style={draw=gray!40,line width=0.4pt,dash pattern=on 1pt off 1.5pt},
  bp key/.style={font=\scriptsize,anchor=west,inner sep=0pt},
  bp use/.style={draw=gray!60,line width=0.5pt,-{Stealth[length=2.6pt,width=2.2pt]}},
  bp stubuse/.style={draw=gray!45,line width=0.45pt,dash pattern=on 1.2pt off 1pt,-{Stealth[length=2.6pt,width=2.2pt]}},
  bp asmuse/.style={draw=orange!65!black,line width=0.5pt,dash pattern=on 2pt off 1.2pt,-{Stealth[length=2.6pt,width=2.2pt]}},
  bp neck/.style={draw=red!60!black,line width=1.1pt,-{Stealth[length=3.6pt,width=3.2pt]}},
  bp long/.style={opacity=0.5}]
  \node[bp unit] at (3.186,3.069) {第~6~章};
  \node[bp unit] at (8.725,3.069) {第~8~章};
  \draw[bp sep] (0.605,2.929) -- (0.605,-2.799);
  \draw[bp sep] (5.941,2.929) -- (5.941,-2.799);
  \draw[bp sep] (11.277,2.929) -- (11.277,-2.799);
  \node[bp R] (n1) at (1.509,0.332) {\detokenize{ee-R-ImPi}};
  \node[bp R] (n2) at (1.509,1.103) {\detokenize{ee-single-}\\\detokenize{scale}};
  \node[bp R] (n3) at (1.509,-0.318) {\detokenize{ee-euclid-ir}};
  \node[bp R] (n4) at (3.099,1.092) {\detokenize{ee-adler-rg}};
  \node[bp R] (n5) at (4.863,1.426) {\detokenize{ee-R-finite}};
  \node[bp R] (n6) at (3.099,-0.858) {\detokenize{ee-oneloop}};
  \node[bp R] (n7) at (1.509,-1.221) {\detokenize{ee-twoloop}};
  \node[bp R] (n8) at (4.863,-1.415) {\detokenize{ee-cutting-}\\\detokenize{exact}};
  \node[bp A] (n9) at (4.863,0.776) {\detokenize{ee-cuts}};
  \node[bp R] (n10) at (3.099,-0.208) {\detokenize{ee-contour}};
  \node[bp R] (n11) at (4.863,-0.644) {\detokenize{ee-R-alphas}};
  \node[bp A] (n12) at (3.099,-1.508) {\detokenize{ee-ope}};
  \node[bp conj] (n13) at (4.863,-2.245) {\detokenize{ee-np-ope}};
  \node[bp R] (n14) at (3.099,0.442) {\detokenize{ee-smearing}};
  \node[bp conj] (n15) at (4.863,0.066) {\detokenize{ee-duality}};
  \node[bp R] (n16) at (7.019,0.695) {\detokenize{pdf-}\\\detokenize{invariance}};
  \node[bp R] (n17) at (7.019,-0.196) {\detokenize{pdf-}\\\detokenize{support-pt}};
  \node[bp A] (n18) at (8.841,-0.013) {\detokenize{pdf-support-}\\\detokenize{hadron}};
  \node[bp R] (n19) at (7.019,-1.087) {\detokenize{pdf-wilson-}\\\detokenize{moments}};
  \node[bp R] (n20) at (8.841,-0.904) {\detokenize{pdf-moments-}\\\detokenize{local}};
  \node[bp R] (n21) at (10.431,-1.127) {\detokenize{pdf-dglap}};
  \node[bp R] (n22) at (10.431,2.428) {\detokenize{pdf-lo-}\\\detokenize{kernels}};
  \node[bp R] (n23) at (10.431,-0.357) {\detokenize{pdf-}\\\detokenize{partonic-dr}};
  \node[bp R] (n24) at (10.431,0.535) {\detokenize{pdf-dis-}\\\detokenize{xspace}};
  \node[bp R] (n25) at (10.431,1.426) {\detokenize{pdf-sum-}\\\detokenize{rules}};
  \node[bp conj] (n26) at (7.019,-2.075) {\detokenize{pdf-}\\\detokenize{positivity}};
  \node[bp A] (n27) at (10.431,-2.305) {\detokenize{pdf-}\\\detokenize{timelike-}\\\detokenize{dglap}};
  \node[bp stub] (n28) at (0.000,-0.035) {第 3 章};
  \node[bp stub] (n29) at (0.000,-0.858) {第 5 章};
  \node[bp stub] (n30) at (11.825,-1.317) {附录 C};
  \node[bp stub] (n31) at (11.825,0.541) {附录 K};
  \node[bp stub] (n32) at (11.825,-0.667) {附录 L};
  \begin{scope}[on background layer]
  \draw[bp stubuse,bp long] (n30) -- (n6);
  \draw[bp stubuse,bp long] (n30) -- (n7);
  \draw[bp stubuse] (n28) -- (n1);
  \draw[bp stubuse] (n28) -- (n3);
  \draw[bp stubuse] (n28) -- (n6);
  \draw[bp stubuse,bp long] (n28) -- (n9);
  \draw[bp stubuse,bp long] (n28) -- (n16);
  \draw[bp stubuse,bp long] (n28) -- (n19);
  \draw[bp stubuse,bp long] (n28) -- (n23);
  \draw[bp stubuse,bp long] (n28) -- (n24);
  \draw[bp stubuse] (n31) -- (n24);
  \draw[bp stubuse,bp long] (n32) -- (n19);
  \draw[bp stubuse] (n32) -- (n20);
  \draw[bp stubuse] (n29) -- (n6);
  \draw[bp use] (n1) -- (n5);
  \draw[bp use] (n1) -- (n6);
  \draw[bp use] (n1) -- (n10);
  \draw[bp use] (n1) -- (n14);
  \draw[bp use] (n2) -- (n4);
  \draw[bp use] (n3) -- (n4);
  \draw[bp use] (n3) -- (n12);
  \draw[bp use] (n4) -- (n5);
  \draw[bp use] (n5) to[bend left=35] (n9);
  \draw[bp use] (n5) to[bend left=35] (n11);
  \draw[bp use] (n6) -- (n8);
  \draw[bp use] (n6) -- (n11);
  \draw[bp use] (n7) -- (n8);
  \draw[bp use] (n7) -- (n11);
  \draw[bp use] (n7) -- (n12);
  \draw[bp use] (n8) to[bend left=35] (n9);
  \draw[bp use] (n10) -- (n11);
  \draw[bp use] (n12) -- (n13);
  \draw[bp use] (n13) to[bend left=35] (n15);
  \draw[bp use] (n14) -- (n15);
  \draw[bp use] (n16) -- (n23);
  \draw[bp use] (n16) -- (n25);
  \draw[bp use] (n17) -- (n18);
  \draw[bp use] (n17) -- (n20);
  \draw[bp use] (n19) -- (n20);
  \draw[bp use] (n20) -- (n21);
  \draw[bp use] (n21) to[bend left=35] (n22);
  \draw[bp use] (n22) to[bend left=35] (n23);
  \draw[bp use] (n22) to[bend left=35] (n25);
  \draw[bp use] (n22) to[bend left=35] (n27);
  \draw[bp use] (n23) to[bend left=35] (n24);
  \end{scope}
  \begin{scope}[shift={(-0.455,-3.299)}]
    \node[bp R,anchor=west,minimum width=4.5mm] (k1) at (0,0) {R};
    \node[bp R,bp cond,anchor=west,minimum width=4.5mm,right=0.8mm of k1] (k2) {R};
    \node[bp A,anchor=west,minimum width=4.5mm,right=0.8mm of k2] (k3) {A};
    \node[bp P,anchor=west,minimum width=4.5mm,right=0.8mm of k3] (k4) {P};
    \node[bp C,anchor=west,minimum width=4.5mm,right=0.8mm of k4] (k5) {C};
    \node[bp V,anchor=west,minimum width=4.5mm,right=0.8mm of k5] (k6) {V};
    \node[bp key,right=1.2mm of k6] (t1) {严格性标签（虚线框为 $R_{\rm cond}$）};
    \node[bp conj,anchor=west,minimum width=7mm,minimum height=3.4mm,right=2.5mm of t1] (k7) {};
    \node[bp key,right=1mm of k7] (t2) {猜想};
    \node[bp asm,anchor=west,minimum width=4.5mm,right=2.5mm of t2] (k8) {};
    \node[bp key,right=1mm of k8] (t3) {编号假设};
    \node[bp math,anchor=west,minimum width=4.5mm,right=2.5mm of t3] (k9) {};
    \node[bp key,right=1mm of k9] (t4) {纯数学};
    \node[bp stub,anchor=west,minimum width=4.5mm,right=2.5mm of t4] (k10) {};
    \node[bp key,right=1mm of k10] (t5) {其他单元};
    \node[bp fold,anchor=west,minimum width=4.5mm,right=2.5mm of t5] (k11) {$+k$};
    \node[bp key,right=1mm of k11] {折叠};
    \draw[bp use] (0,-0.42) -- ++(6mm,0) coordinate (e1);
    \node[bp key] (u1) at ($(e1)+(1.2mm,0)$) {依赖（前提指向陈述）};
    \draw[bp asmuse] ($(u1.east)+(2.5mm,0)$) -- ++(6mm,0) coordinate (e2);
    \node[bp key] (u2) at ($(e2)+(1.2mm,0)$) {编号假设 $\to R_{\rm cond}$};
    \draw[bp neck] ($(u2.east)+(2.5mm,0)$) -- ++(6mm,0) coordinate (e3);
    \node[bp key] (u3) at ($(e3)+(1.2mm,0)$) {瓶颈依赖（主表）};
    \draw[bp use,bp long] ($(u3.east)+(2.5mm,0)$) -- ++(6mm,0) coordinate (e4);
    \node[bp key] (u4) at ($(e4)+(1.2mm,0)$) {跨两列以上};
    \draw[bp stubuse] ($(u4.east)+(2.5mm,0)$) -- ++(6mm,0) coordinate (e5);
    \node[bp key] at ($(e5)+(1.2mm,0)$) {来自其他单元};
  \end{scope}
\end{tikzpicture}
\caption{第二部分（第~6--8~章）主要定理的依赖图：总表中的陈述与本部分的全部定理、它们在本部分内的直接前提及所用的编号假设（27 个陈述节点、5 个汇总节点）。图例与画法同\cref{fig:appP-dependency}；灰色框代表本部分之外的单元，“$+k$”为按折叠规则省略的 $k$ 个次要陈述。}\label{fig:appP-dependency-c}
\end{figure}
